\par\Needspace{7\baselineskip}\bigskip
\begin{center}
\SourceLeaf{442}
\pdfbookmark[2]{Gauss an Hansen, Göttingen, 11. März 1843}{pallas-brief-18430311}
\textsc{Gauss} an \textsc{Hansen}, Göttingen, 11. März 1843.
\end{center}
\noindent\hspace*{\parindent}Von einer Woche zur andern ist mein Dank für die gefällige Übersendung des Berliner Monatsberichts (der um dieselbe Zeit auch auf gewöhnlichem Wege mir zu Gesicht kam) verschoben, weil ich hoffte einige Zeit zu gewinnen, in den Gegenstand etwas weiter in meiner Antwort eingehen zu können. Leider ist diese Hoffnung getäuscht, und selbst in den bevorstehenden Ferien, für welche sich schon im Voraus so viele Rückstände und neue Abhaltungen gesammelt haben, darf ich mir kaum Hoffnung für einige freie Musse machen.
\par Ich beschränke mich daher darauf, meine Freude darüber auszusprechen, dass Sie bei den Perturbationsrechnungen auf ähnliche Art verfahren, wie ich schon vor mehr als 30 Jahren bei meinen weitumfassenden Rechnungen über die Pallasstörungen zu Werke gegangen bin, in so fern Sie den Gebrauch von Reihen nach den Excentricitäten und Neigungen ganz cassiren. Freilich haben Sie für die Cometenstörungen auch ganz besondere noch andere Untersuchungen nöthig gehabt, zu denen für die Pallas keine Veranlassung sich fand und überhaupt in vielen andern Beziehungen abweichende Wege eingeschlagen. Bei den Störungen der Pallas durch Jupiter brauchte ich die Methode der variabeln Elemente und zwar mit Vorbedacht, denn obgleich man so 6 Elemente zu behandeln hat, während bei dem andern Verfahren nur halb so viele sind, so habe ich doch für den praktischen Gebrauch jenes vorgezogen; man braucht die Rechnung für jedes Jahr nur einmal (für Einen Tag) zu machen, wozu ich eine besondere Hülfstafel construirte, vermittelst welcher in vergleichungsweise sehr kurzer Zeit die Rechnung absolvirt werden kann, obgleich zusammen 801 Gleichungen (1602, wenn die Sinus und Cosinus-Glieder desselben Arguments getrennt gezählt werden; es sind vollständig alle, deren Coëfficient über \(0''{,}1\) geht) berücksichtigt worden\footnote{Es ist wohl unnöthig zu bemerken, dass die Arbeit, von der ich jetzt spreche, \textls[180]{ganz} verschieden ist von der gleichfalls von mir für den ganzen Zeitraum von 1802 bis etwa 1818 oder 1820 fortgeführten Rechnung durch Quadratur; diese nenne ich specielle, jene generelle Rechnung und letztere hat in jener eine bei so ausgedehnten Rechnungen höchst nothwendige Controlle gefunden.}. Durch sehr einfache Mittel kann man dann das Resultat für die ganze \SourceLeaf{443}Beobachtungs-Saison ausreichend machen, ohne der Schärfe etwas zu vergeben. Auch die Störung durch Saturn wurde berechnet und die durch Mars nach einer wesentlich verschiedenen Methode angefangen, aber nicht vollendet. Andere immer weiter sich verzweigende Geschäfte haben mir später gar nicht erlaubt, auf jene Arbeiten wieder zurückzukommen, und es steht dahin, ob in meinem Alter ich Musse und Lust haben werde, mich wieder in die Sachen hineinzuarbeiten, da noch so viele andere Dinge sind, die ich eben so ungern untergehen lassen möchte. Sie sind sehr glücklich, dass Sie in einer äussern Lage sind, wo Sie Ihre Zeit nicht zu versplittern brauchen, und es wird mir jedenfalls zur Beruhigung gereichen, dass dieser Zweig der Astronomie bei Ihnen in den besten Händen ist. Ob ich bei den Marsstörungen etwas mit Ihrer mir angezeigten Integrationsart zusammenhangendes gebraucht habe, kann ich jetzt, wo mir die Sachen seit fast 25 Jahren entfremdet sind, nicht bestimmt ermitteln, möchte es aber fast bezweifeln.\quad\mbox{. . . . .}
\par
